\section{Aims and scope}

\noindent\emph{Status: draft -- the purpose, the division of labour between the sections and the notation are
settled; the framing will be adjusted as milestones land.}

\texttt{emfit} is meant to be an engineer's tool rather than a demonstration of a method. Given a
three-dimensional structure -- a cavity, a piece of laminate, a patch of a printed circuit board -- together with
a description of the materials filling it, the tool should answer the questions an engineer actually asks: where
the field is, at which frequencies the structure rings, and how much power leaves through which port. The
discretisation used throughout is the finite integration technique (FIT) introduced by Weiland
\cite{weiland1977,weiland1996} and systematised, in the form adopted here, by Clemens and Weiland
\cite{clemens2001} and Schuhmann and Weiland \cite{schuhmann2001}. The rest of this report is about how far that
single sentence has been made true, and where it has not.

\subsection{What the tool is for}

Two high-frequency workflows set the direction of the work, and a third, low-frequency extension follows once
they are solid.

\begin{itemize}[leftmargin=1.4em,itemsep=2pt]
  \item \textbf{A general three-dimensional solver.} A box of (for now isotropic, lossless) material is excited by
        a soft current source, and the fields are stepped in time. What is wanted out of it is the ring-down
        spectrum of a cavity, field maps on a cut plane, and an energy balance that can be trusted because it is
        exact rather than merely small.
  \item \textbf{An antenna on a printed circuit board.} A driven structure with feed and load ports, from which
        scattering parameters and a radiation pattern are extracted. This is the commercially relevant case and
        it is the reason the roadmap insists on ports, on a near-to-far-field transform and on importing real
        geometry, rather than on ever larger toy cavities; the antenna-theory side of the acceptance tests
        follows \cite{balanis2016}.
  \item \textbf{A low-frequency extension.} Magnetoquasistatic eddy currents in laminations and conductors, and a
        time-parallel integrator for long transient runs. This is deliberately last: it shares the grid and the
        material layer, and nothing else.
\end{itemize}

\subsection{Why the finite integration technique}

The choice of method is not incidental; six properties carried the decision.

FIT discretises the \emph{integral} form of Maxwell's equations rather than the differential form. The unknowns
are line integrals of the electric and magnetic field along primal and dual edges, and surface integrals of flux
through primal and dual faces \cite{weiland1977,clemens2001}. Because the integral form is an exact statement
about any cell, not only about an infinitesimally small one, there is no Taylor expansion anywhere in the
derivation: the only step at which the discrete model departs from Maxwell's equations is the constitutive step,
where the material response is replaced by a matrix on the grid.

Topology is exact and \emph{metric free}. The discrete operators are incidence matrices whose entries are $0$ and
$\pm 1$ and nothing else, so the identities $\Cmat\Gmat = 0$ and $\Smat\Cmat = 0$ hold in integer arithmetic on
every grid, uniform or not. Refining one region of the domain therefore changes nothing but the metric, i.e.\ the
Hodge star; that is a property a general finite-element assembly does not give for free, and it is what makes
non-uniform meshing cheap here.

The scheme is energy conserving in a strong sense, made precise in the section on energy conservation and stability: the leapfrog update
has an exact discrete invariant, a mixed-state energy that includes a first-order cross term, and that invariant
is conserved at the level of round-off rather than merely bounded.

There is no linear system to solve in the time domain. The Hodge stars are diagonal, so a step is two sparse
matrix-vector products: the electric and magnetic unknowns update each other explicitly. No global matrix is
factored, no preconditioner is built, and no iterative tolerance has to be chosen -- which is why a step is cheap
enough that the time domain is a practical way to obtain a spectrum.

FIT is the parent of Yee's scheme \cite{yee1966}. On a uniform orthogonal grid with staggered unknowns the two
coincide, so fifty years of accumulated experience with finite-difference time-domain methods on dispersion,
stability and staircasing \cite{taflove2005} transfers directly, and the FIT derivation supplies the reason why
those properties hold.

The price is regularity. FIT lives on structured hexahedral grids, so a curved surface becomes a staircase.
Finite elements on unstructured tetrahedra, or edge elements in the sense of Bossavit \cite{bossavit1998} and
Whitney forms in the sense of \cite{jin2014}, fit geometry better but pay for it with a global sparse solve at
every step, or with an implicit formulation and its conditioning problems. For the class of problems this project
targets -- boxes, cavities, layered laminates, antennas over a finite ground plane -- the structured grid is not
a practical limitation, and the cheap explicit step is decisive.

\subsection{Scope of this report}

This is a living document. It records what is implemented, the mathematics the implementation rests on, and the
measured evidence for each claim; every number quoted in the validation section is asserted by a test in the
repository, so a claim that stops being true turns the suite red. Anything that is designed but not implemented is
confined to the Outlook and to \texttt{docs/roadmap.md}; it is not described here in the present
tense. Where the two disagree, the roadmap is authoritative for status and this report for method.

\subsection{Implemented and not yet implemented}

\begin{itemize}[leftmargin=1.4em,itemsep=2pt]
  \item[] \textbf{Working today.}
        Orthogonal grids with non-uniform spacing and their dual cell complexes (\texttt{emfit/grid.py},
        \texttt{emfit/operators.py}); the three topological operators with $\Cmat\Gmat = 0$ and $\Smat\Cmat = 0$
        exact, verified against hand-checked stencils; the diagonal Hodge stars $\Em$, $\Mm$ and $\Ms$; the
        assembled system \texttt{FITSystem} with PEC walls imposed as a constraint on the unknown vector; the
        explicit leapfrog stepper, whose mixed-state invariant drifts by $8.5\times10^{-16}$ relative over the
        $343$~ns ring-down of the example (1800 steps) and stays at the $10^{-15}$ level over 2000 steps of random
        fields; the sharp CFL bound, stable at $1.01$ times the limit and divergent at $1.02$; a soft current
        source with the power balance $P = \vec e\cdot\vec j$; a dense cavity eigensolver with static-mode
        filtering and an analytic rectangular-box reference; and a headless visualisation layer with a runnable
        example. The suite is $336$ tests in $24$ modules; the fast selection runs $327$ of them in about
        $17$~s and deselects $9$ slow ones.
  \item[] \textbf{Measured accuracy.} On a $1.0\times0.7\times0.5$~m PEC box the first six modes are reproduced
        with a mean absolute error of $1.213\,\%$ and a maximum of $1.390\,\%$ on the $10\times7\times5$ grid, and
        $0.615\,\%$ / $0.711\,\%$ on $14\times10\times7$ -- a mean-error ratio of $0.507$ against the $0.510$
        predicted by second order for the $1.4$ times refinement used.
  \item[] \textbf{Absorbing boundaries, in part.} The first-order condition of G.~Mur \cite{mur1981} is
        implemented and measured (Section~\ref{sec:res-abc}): an outgoing pulse is absorbed to the level the
        discretisation allows, so a domain need no longer be closed by PEC or PMC walls. A matched lossy layer
        stands in for the perfectly matched layer of \cite{berenger1994} and is implemented too
        (Section~\ref{sec:res-pml}): it absorbs far more deeply than the condition does on the same geometry,
        but a layer of finite-difference cells is not a perfect match, and the acceptance criterion attached to
        it is not met. The near-to-far-field transform is implemented and checked against a closed-form source,
        and a bistatic scattering width computed through it has been set beside a Mie series
        (Section~\ref{sec:res-scatter}), though the agreement the milestone asks for has not been reached.
        Waveguide ports and S-parameters remain unavailable, and with them every input impedance.
  \item[] \textbf{Not implemented.} Geometry import from a board layout. A large-grid eigensolver, the dense
        path refusing more than 3000 free edges. Magnetoquasistatic eddy currents and time-parallel
        integration. Lossy and dispersive media: $\Ms$ is assembled but not yet coupled into the step, and
        the material layer understands diagonal tensors only. There is no non-orthogonal grid and no
        finite-element path; the two-dimensional studies of Section~\ref{sec:results} are the transverse-magnetic
        case of a grid one cell thick in the normal direction, which is what a single cell gives.
\end{itemize}

\subsection{Notation and conventions}

\begin{table}[H]
  \centering
  \caption{Symbols used throughout. The unit is the unit of the corresponding discrete quantity, which is what
  the code stores; the continuum field is the integrand, not the unknown.}
  \label{tab:notation}
  \begin{tabular}{@{}lll@{}}
    \toprule
    Symbol & Meaning & Unit \\
    \midrule
    $\vec e$      & electric field strength            & V/m   \\
    $\vec h$      & magnetic field strength            & A/m   \\
    $\vec b$      & magnetic flux density              & Wb/m$^2$ \\
    $\vec d$      & electric flux density              & C/m$^2$ \\
    $\vec j$      & electric current density           & A/m$^2$ \\
    $\Em$         & diagonal Hodge star, permittivity  & F     \\
    $\Mm$         & diagonal Hodge star, permeability  & H     \\
    $\Cmat$       & primal curl (faces $\times$ edges) & 1     \\
    $\Gmat$       & primal gradient (edges $\times$ nodes) & 1 \\
    $\Smat$       & primal divergence (cells $\times$ faces) & 1 \\
    $\eps_0$      & vacuum permittivity                & F/m   \\
    $\mu_0$       & vacuum permeability                & H/m   \\
    $c_0$         & vacuum speed of light              & m/s   \\
    \bottomrule
  \end{tabular}
\end{table}

\begin{itemize}[leftmargin=1.4em,itemsep=2pt]
  \item SI units throughout, with no unitless rescaling anywhere in the code: lengths in metres, energies in
        joules, powers in watts. All internal quantities are \texttt{float64}.
  \item Real-valued time-domain fields everywhere except in the eigenmode analysis, where complex phasors are used; the phasor convention is stated there where it
        matters and nowhere else.
  \item Degrees of freedom are written as bold lower case letters, and a tilde marks a quantity on the dual
        complex: $\mathbf{e}$, $\mathbf{h}$, and $\tilde{\mathbf{h}}$ for the dual line integral of $\vec h$.
  \item Unknown vectors are the concatenation \texttt{EX}, \texttt{EY}, \texttt{EZ} of the primal edges, then
        \texttt{FX}, \texttt{FY}, \texttt{FZ} of the primal faces, with each block flattened in C order; material
        vectors are indexed exactly like the block they belong to, not like the cells.
  \item Three dimensions only. There is no two-dimensional transverse-magnetic reduction, and every statement in
        this report is about a hexahedral complex.
\end{itemize}
